%We focus on the general task of mapping one variable-length sequence of symbol representations ${x_1, ..., x_n} \in \mathbb{R}^d$ to another sequence of the same length ${y_1, ..., y_n} \in \mathbb{R}^d$. \marginpar{should we use this notation? alternatively we can just say "d-dimensional vectors"}

在本节中，我们比较自注意力层与循环和卷积层的各个方面，这些层通常用于将一个可变长度的符号表征序列 $(x_1, ..., x_n)$ 映射到另一个等长序列 $(z_1, ..., z_n)$，其中 $x_i, z_i \in \mathbb{R}^d$，例如在典型的序列转导编码器或解码器中的隐层。为了支持我们对自注意力的使用，我们考虑三个期望。

一个是每层的总计算复杂度。
另一个是能够并行化的计算量，通过所需的最少顺序操作数来衡量。

第三个是网络中长程依赖之间的路径长度。学习长程依赖是许多序列转导任务中的关键挑战。影响学习这样的依赖能力的一个关键因素是前向和反向信号必须在网络中遍历的路径长度。输入和输出序列中任何位置组合之间的这些路径越短，学习长程依赖就越容易 \citep{hochreiter2001gradient}。因此，我们还比较由不同层类型组成的网络中任意两个输入和输出位置之间的最大路径长度。

%\subsection{Computational Performance and Path Lengths}

\begin{table}[t]
\caption{
  不同层类型的最大路径长度、每层复杂度和最少顺序操作数。$n$ 是序列长度，$d$ 是表征维度，$k$ 是卷积的核大小，$r$ 是受限自注意力中邻域的大小。}
  %Attention models are quite efficient for cross-positional communications when sequence length is smaller than channel depth.
\label{tab:op_complexities}
\begin{center}
\vspace{-1mm}
%\scalebox{0.75}{

\begin{tabular}{lccc}
\toprule
Layer Type & Complexity per Layer & Sequential & Maximum Path Length  \\
           &             & Operations &   \\
\hline
\rule{0pt}{2.0ex}Self-Attention & $O(n^2 \cdot d)$ & $O(1)$ & $O(1)$ \\
Recurrent & $O(n \cdot d^2)$ & $O(n)$ & $O(n)$ \\

Convolutional & $O(k \cdot n \cdot d^2)$ & $O(1)$ & $O(log_k(n))$ \\
%\cmidrule
Self-Attention (restricted)& $O(r \cdot n \cdot d)$ & $O(1)$ & $O(n/r)$ \\

%Convolutional (separable) & $O(k \cdot n \cdot d + n \cdot d^2)$ & $O(1)$ & $O(log_k(n))$ \\

%Position-wise Feed-Forward & $O(n \cdot d^2)$ & $O(1)$ & $\infty$ \\

%Fully Connected & $O(n^2 \cdot d^2)$ & $O(1)$ & $O(1)$ \\
%Convolutional (separable) & $O(k \cdot n \cdot d + n \cdot d^2)$ & $O(1)$ & $O(log_k(n))$ \\

%Position-wise Feed-Forward & $O(n \cdot d^2)$ & $O(1)$ & $\infty$ \\

%Fully Connected & $O(n^2 \cdot d^2)$ & $O(1)$ & $O(1)$ \\
\bottomrule
\end{tabular}
%}
\end{center}
\end{table}


%\begin{table}[b]
%\caption{
%  Maximum path lengths, per-layer complexity and minimum number of sequential operations for different layer types. $n$ is the sequence length, $d$ is the representation dimensionality, $k$ is the kernel size of convolutions and $r$ the size of the neighborhood in localized self-attention.}
  %Attention models are quite efficient for cross-positional communications when sequence length is smaller than channel depth.
%\label{tab:op_complexities}
%\begin{center}
%\vspace{-1mm}
%%\scalebox{0.75}{
%
%\begin{tabular}{lccc}
%\hline
%Layer Type & Receptive & Complexity per Layer & Sequential  %\\
%           & Field Size     &            & Operations  \\
%\hline
%Self-Attention & $n$ & $O(n^2 \cdot d)$ & $O(1)$ \\
%Recurrent & $n$ & $O(n \cdot d^2)$ & $O(n)$ \\

%Convolutional & $k$ & $O(k \cdot n \cdot d^2)$ & %$O(log_k(n))$ \\
%\hline 
%Self-Attention (localized)& $r$ & $O(r \cdot n \cdot d)$ & %$O(1)$ \\

%Convolutional (separable) & $k$ & $O(k \cdot n \cdot d + n %\cdot d^2)$ & $O(log_k(n))$ \\

%Position-wise Feed-Forward & $1$ & $O(n \cdot d^2)$ & $O(1)$ %\\

%Fully Connected & $n$ & $O(n^2 \cdot d^2)$ & $O(1)$ \\

%\end{tabular}
%%}
%\end{center}
%\end{table}

%The receptive field size of a layer is the number of different input representations that can influence any particular output representation. Recurrent layers and self-attention layers have a full receptive field equal to the sequence length $n$. Convolutional layers have a limited receptive field equal to their kernel width $k$, which is generally chosen to be small in order to limit computational cost.

如表 \ref{tab:op_complexities} 所述，自注意力层用常数个顺序执行的操作连接所有位置，而循环层需要 $O(n)$ 个顺序操作。
在计算复杂度方面，当序列长度 $n$ 小于表示维度 $d$ 时，自注意力层比循环层快，这对于最先进的机器翻译模型中使用的句子表示（如词片段 \citep{wu2016google} 和字节对 \citep{sennrich2015neural} 表示）来说通常是这样。
为了改进涉及很长序列的任务的计算性能，自注意力可以限制为仅考虑在输入序列中以各自输出位置为中心的大小为 $r$ 的邻域。这会将最大路径长度增加到 $O(n/r)$。我们计划在未来工作中进一步调查这种方法。

核宽度为 $k < n$ 的单个卷积层不连接所有的输入和输出位置对。要做到这一点，在连续核的情况下需要 $O(n/k)$ 个卷积层的堆栈，或在膨胀卷积的情况下需要 $O(log_k(n))$ \citep{NalBytenet2017}，增加了网络中任意两个位置之间的最长路径长度。
卷积层通常比循环层贵得多，因子为 $k$。然而，可分离卷积 \citep{xception2016} 大大降低了复杂度，降至 $O(k \cdot n \cdot d + n \cdot d^2)$。但即使 $k=n$，可分离卷积的复杂度等于自注意力层和逐点前馈层的组合，这是我们模型中采用的方法。

%\subsection{Unfiltered Bottleneck Argument}

%An orthogonal argument can be made for self-attention layers based on when the layer imposes the bottleneck of mapping all of the information used to compute a given output position into a single, fixed-length vector. ...

%There is a second argument for self-attention layers which we call the unfiltered bottleneck argument.   In both recurrent and the convolutional layers, the information that position $i$ receives from the other positions is compressed to a vector of dimension $d$ before it ever can be filtered by the content $x_i$.  More precisely, we can express $y_i = F(i, x_i, G(i, \{x_{j \neq i}\}))$, where $G(i, \{x_{j \neq i}\})$ is a vector of dimension $d$.  Intuitively, we would expect that this would cause a large amount of irrelevant information to crowd out the relevant information.  Self-attention does not suffer from the unfiltered bottleneck problem, since the aggregation happens after filtering, and so, intuitively, we have the chance of transmitting lots of relevant information.

作为附带好处，自注意力可以产生更可解释的模型。我们检查了模型的注意力分布，并在附录中呈现和讨论了示例。不仅个别注意力头清楚地学会执行不同的任务，许多还表现出与句子句法和语义结构相关的行为。

